\begin{table}[htbp]
\centering
\caption{A fixed detection-capacity budget, every fourth grid point of 21.}
\label{tab:budget}
\small
\begin{tabular}{ccccccccc}
\toprule
$k$ & $d_0$ & $d_1$ & \shortstack{average\\detection} & \shortstack{recorded\\rate}
& AUC vs $F$ & AUC vs $C$ & Brier vs $F$ & Brier vs $C$ \\
\midrule
0.00 & 0.400 & 0.400 & 0.400 & 0.080 & 0.5000 & 0.5000 & 0.0736 & 0.1744 \\
0.16 & 0.360 & 0.520 & 0.400 & 0.080 & 0.5408 & 0.5000 & 0.0734 & 0.1746 \\
0.32 & 0.320 & 0.640 & 0.400 & 0.080 & 0.5815 & 0.5000 & 0.0728 & 0.1752 \\
0.48 & 0.280 & 0.760 & 0.400 & 0.080 & 0.6223 & 0.5000 & 0.0719 & 0.1761 \\
0.64 & 0.240 & 0.880 & 0.400 & 0.080 & 0.6630 & 0.5000 & 0.0705 & 0.1775 \\
0.80 & 0.200 & 1.000 & 0.400 & 0.080 & 0.7038 & 0.5000 & 0.0688 & 0.1792 \\
\bottomrule
\end{tabular}
\tablenotes{What is held fixed is $\bar d$, the average probability of detecting an
existing deficiency; it is not a measure of audit hours or cost.
Detection is set by $d_0 = \bar d - \pi k$ and $d_1 = \bar d + (1-\pi)k$, with
$p = 0.20$, $\pi = 0.25$ and
$\bar d = 0.40$, all prespecified. Average detection and
recorded prevalence are invariant in $k$ by construction; the table shows they are.
Exact calculation.}
\end{table}
